\scanpage{123}
\mspage{61}
en ces paramètres, et en se ramenant à un $f,g$ premiers entre eux, \unsure{on est|on voit} --- faisant encore éclater $a$ --- que les diviseurs de $f,g$ \deleted{\unsure{sont|se}} \unsure{séparent|séparés}\ldots

b) Un anneau de val $V$ \deleted{qui} pour l'extension $L/K$ domine un pt bien déterminé de chaque $X$, d'où un $x\in\mathfrak X=\varprojlim X$, et $V$ domine $\underline{O}_{\mathfrak X,x}$, donc lui est égal puisque ce dernier est un anneau de valuation.

\uline{NB} Tout \inserted{modèle $U$ de $L/K$ définit un} \uline{ouvert} de $\mathfrak X$ \deleted{\unsure{correspondant}} \deleted{\unsure{évidemment}} \inserted{(en le plongeant} dans un modèle propre $X$, c'est clair) et \deleted{et} les modèles correspondent à \uline{certains} ouverts de $\mathfrak X$ (\uline{pas} stables par réunions finies).

Pour \unsure{tt|$\forall$} anneau \unsure{loc\textsuperscript{l}|local} \inserted{\unsure{norm\textsuperscript{l}|normal}} $O\subset L$ de corps des fractions $L$, on peut regarder $\operatorname{Spec}\underline{O}=S$, son pt fermé $\unsure{x|s}$, le \unsure{normali-}
\margin{122}
